%\subsection{Aufgaben}
%	Laplace-Transformierte berechnen, zum Beispiel f�r folgende Manipulationen:
%\begin{itemize}
%	\item Transformation des Integrals zeigen
%	\item Differentiation der Bildfuktion zeigen
%	\item Integration des Bildfunktion
%	\item D�mpfung
%	\item Periodische Funktionen
%\end{itemize}

\ifintegriereaufgaben

\section*{Laplace Transformation Aufgaben}\label{MA2-laplacetransformation-aufgaben}
Erkl�rungen und Beispiele, die zur Bearbeitung der Aufgaben hilfreich sind, finden Sie im Skript\wennreferenzvorhanden{MA2-laplacetransformation}{ (Abschnitt \abschnittsverweis{MA2-laplacetransformation})} und, falls Sie nicht weiterkommen, dann schauen Sie \hyperref[MA2-laplacetransformation-loesungen]{hier}.
\else
\section{Laplace Transformation Aufgaben}\label{MA2-laplacetransformation-aufgaben}

Regeln und Beispiele zu diesen Aufgaben finden Sie im Skript. Die L�sungen finden Sie \hyperref[MA2-laplacetransformation-loesungen]{hier}.
\fi

\begin{exercisebox}[Uneigentliche Integrale] %A.Z. Integralrechnung  % Papula I.V, Abschnitt 9, Aufgabe 1 und 2
	Bestimmen Sie den Wert der folgenden uneigentlichen Integrale, falls er existiert.
	\begin{eqnarray*}
		(a)\;\int_0^\infty e^{-x}\,\mathrm{d}x\quad
		(b)\;\int_0^\infty x\cdot e^{-x}\,\mathrm{d}x\quad
		(c)\;\int_{-\infty}^2 e^{x}\,\mathrm{d}x\quad
		(d)\;\int_{0}^\infty x^2\cdot e^{-2x}\,\mathrm{d}x
	\end{eqnarray*}
\end{exercisebox}

\begin{exercisebox}[Uneigentliche Integrale] %A.Z. Integralrechnung  % Papula I.V, Abschnitt 9, Aufgabe 4
	Bestimmen Sie den Wert der folgenden uneigentlichen Integrale, falls er existiert.
	\begin{eqnarray*}
		(a)\;\int_{-1}^0 \frac{1}{\sqrt{1+x}}\,\mathrm{d}x\quad
		(b)\;\int_{-1}^{1} \frac{1}{x^2}\,\mathrm{d}x\quad
		(c)\;\int_{0}^{10} \frac{e^{x}}{e^x-1}\,\mathrm{d}x
	\end{eqnarray*}
\end{exercisebox}


\begin{exercisebox}[Laplace-Transformation] % Papula 2 VI, Abschnitt 1, Aufgabe 1 und 3
	Bestimmen Sie f�r folgende Funktionen die Laplace-Transformierte durch explizite Berechnung.
	\begin{itemize}
	\item[(a)] $f(t) = 2t\cdot e^{-4t}$
	\item[(b)] $f(t) = e^{-\delta t} \cdot \sin(\omega t), \; \delta, \omega >0$
	\item[(c)] $f(t) = \mathrm{sinh}(at),\, a>0$
	\item[(d)] $f(t) = \sin^2(t)$
\item[(e)] (Rechteckimpuls) F�r $a>0$
			\[ f(t) = \begin{cases} A, \quad & 0<t<a\,,\\
					-A, \quad & a< t < 2a\,,\\
			0, \quad \mathrm{sonst}\,. \end{cases}\]
		\item[(f)] (Sinusimpuls) F�r $a>0$
			\[ f(t) = \begin{cases} A\cdot \sin\left( \frac{\pi\cdot t}{a}\right), \quad & 0\leq t \leq a\,,\\
			0, \quad \mathrm{sonst}\,. \end{cases}\]
	\end{itemize}
\end{exercisebox}

\begin{exercisebox}[Laplace-Transformation] % Papula 2, VI, Abschnitt 2, Aufgabe 4
	Bestimmen Sie
	%mit Hilfe von
	%\[ \sin(t) \,\laplace\, \frac{1}{s^2+1} \]
	die Laplace-Transformierte von
	\[ f(t) = \sin(\omega t + \varphi), \; \omega,\varphi >0\,. \]
	%indem Sie zuerst eine Verschiebung und dann eine Skalierung verwenden.
\end{exercisebox}
